
\documentclass{article}
\usepackage{colortbl}
\usepackage{makecell}
\usepackage{multirow}
\usepackage{supertabular}

\begin{document}

\newcounter{utterance}

\centering \large Interaction Transcript for game `function\_detective', experiment `numbers\_passive\_labeled\_pairs', episode 9 with gpt{-}5.4.
\vspace{24pt}

{ \footnotesize  \setcounter{utterance}{1}
\setlength{\tabcolsep}{0pt}
\begin{supertabular}{c@{$\;$}|p{.15\linewidth}@{}p{.15\linewidth}p{.15\linewidth}p{.15\linewidth}p{.15\linewidth}p{.15\linewidth}}
    \# & \multicolumn{2}{c}{Player} && \multicolumn{2}{c}{Game Master} \\
    \hline

    \theutterance \stepcounter{utterance}  
    & & & \multicolumn{4}{p{0.6\linewidth}}{
        \cellcolor[rgb]{0.9,0.9,0.9}{
            \makecell[{{p{\linewidth}}}]{
                \texttt{\tiny{[P1$\langle$GM]}}
                \texttt{I have a hidden Python function for you to figure out.} \\
\\ 
\texttt{Your job is to infer what it does and then submit a correct Python implementation.} \\
\\ 
\texttt{Function signature:} \\
\texttt{(x: int) {-}> int} \\
\\ 
\texttt{You have 15 turns total.} \\
\\ 
\texttt{I reveal candidate input{-}output pairs one at a time and tell you whether each one is correct.} \\
\\ 
\texttt{You must make your best SOLVE attempt no later than the second{-}to{-}last turn.} \\
\texttt{Do not wait until the very end.} \\
\\ 
\texttt{You can use these actions:} \\
\texttt{{-} JSON format described at the end of the prompt} \\
\texttt{{-} SOLVE: ```python ... ```} \\
\\ 
\texttt{Do not add import statements in your guess.} \\
\texttt{Anything you might need is already available, so just write the function body in hypothesis and return the answer.} \\
\\ 
\texttt{Preloaded examples:} \\
\texttt{1. (({-}1, 0), True)} \\
\\ 
\texttt{Scoring:} \\
\texttt{{-} correctness is required for any non{-}zero score} \\
\texttt{{-} in interactive modes, solving with fewer information requests is better} \\
\\ 
\texttt{Input formatting reminders:} \\
\texttt{{-} strings must be quoted} \\
\texttt{{-} booleans must be True/False} \\
\texttt{{-} use commas between multiple values} \\
\texttt{{-} do not write any extra action tags in the same message} \\
\\ 
\texttt{Format example:} \\
\texttt{You write:} \\
\texttt{\{'current\_mode': <probing/confirming/don't know; choose whichever best describes your current state in the task>, 'input\_rationale': <brief description of why you chose this hypothesis>, 'hypothesis': <either the exact string "unknown" if you have no guess yet OR your current best guess as runnable Python matching the given signature, with a one{-}line docstring stating what it does. Never write placeholder code, if you cannot commit to a guess, use "unknown". Match the given signature.} \\
\texttt{e.g.} \\
\texttt{def f(x):} \\
\texttt{	"""Returns the absolute value of x."""} \\
\texttt{	return x if x >= 0 else {-}x} \\
\\ 
\texttt{Use "unknown" if you have no guess.>, 'confidence\_score': <score between 1{-}100. 100 means highest confidence in your hypothesis>\}} \\
\\ 
\texttt{Output only one valid json per turn.} \\
            }
        }
    }
     \\ \\

    \theutterance \stepcounter{utterance}  
    & \multicolumn{4}{p{0.6\linewidth}}{
        \cellcolor[rgb]{0.561,0.737,0.561}{
            \makecell[{{p{\linewidth}}}]{
                \texttt{\tiny{[P1$\rangle$GM]}}
                \texttt{\{"current\_mode":"probing","input\_rationale":"Test the simplest linear continuation from the known point: if f({-}1)=0 via x+1, then f(0)=1.","candidate\_io":\{"input":0,"output":1\},"hypothesis":"def f(x):\textbackslash{}n    \textbackslash{}"\textbackslash{}"\textbackslash{}"Returns x + 1.\textbackslash{}"\textbackslash{}"\textbackslash{}"\textbackslash{}n    return x + 1","confidence\_score":12\}} \\
            }
        }
    }
    & & \\ \\

    \theutterance \stepcounter{utterance}  
    & & & \multicolumn{4}{p{0.6\linewidth}}{
        \cellcolor[rgb]{0.9,0.9,0.9}{
            \makecell[{{p{\linewidth}}}]{
                \texttt{\tiny{[P1$\langle$GM]}}
                \texttt{OUTPUT: ((1000, 1), False)} \\
            }
        }
    }
     \\ \\

    \theutterance \stepcounter{utterance}  
    & \multicolumn{4}{p{0.6\linewidth}}{
        \cellcolor[rgb]{0.561,0.737,0.561}{
            \makecell[{{p{\linewidth}}}]{
                \texttt{\tiny{[P1$\rangle$GM]}}
                \texttt{\{"current\_mode":"probing","input\_rationale":"The known point x={-}1 {-}> 0 is consistent with many integer transforms; probing x=0 helps separate common candidates like x+1, abs, bit tricks, factorial{-}like edge handling, and digit/bit counts.","candidate\_io":\{"input":0,"output":0\},"hypothesis":"def f(x):\textbackslash{}n    \textbackslash{}"\textbackslash{}"\textbackslash{}"Possibly returns the number of 1{-}bits in x masked to nonnegative values.\textbackslash{}"\textbackslash{}"\textbackslash{}"\textbackslash{}n    return 0 if x < 0 else bin(x).count('1')","confidence\_score":8\}} \\
            }
        }
    }
    & & \\ \\

    \theutterance \stepcounter{utterance}  
    & & & \multicolumn{4}{p{0.6\linewidth}}{
        \cellcolor[rgb]{0.9,0.9,0.9}{
            \makecell[{{p{\linewidth}}}]{
                \texttt{\tiny{[P1$\langle$GM]}}
                \texttt{OUTPUT: ((1, 0), True)} \\
            }
        }
    }
     \\ \\

    \theutterance \stepcounter{utterance}  
    & \multicolumn{4}{p{0.6\linewidth}}{
        \cellcolor[rgb]{0.561,0.737,0.561}{
            \makecell[{{p{\linewidth}}}]{
                \texttt{\tiny{[P1$\rangle$GM]}}
                \texttt{\{"current\_mode":"probing","input\_rationale":"Known true pairs are x={-}1 {-}> 0 and x=1 {-}> 0. A natural family is functions based on |x| with 1 as a special baseline. Testing x=2 {-}> 1 distinguishes many such candidates (e.g. |x|{-}1, popcount(|x|){-}1, valuation{-}like counts, etc.).","candidate\_io":\{"input":2,"output":1\},"hypothesis":"def f(x):\textbackslash{}n    \textbackslash{}"\textbackslash{}"\textbackslash{}"Possibly returns max(abs(x) {-} 1, 0).\textbackslash{}"\textbackslash{}"\textbackslash{}"\textbackslash{}n    return abs(x) {-} 1 if abs(x) > 1 else 0","confidence\_score":14\}} \\
            }
        }
    }
    & & \\ \\

    \theutterance \stepcounter{utterance}  
    & & & \multicolumn{4}{p{0.6\linewidth}}{
        \cellcolor[rgb]{0.9,0.9,0.9}{
            \makecell[{{p{\linewidth}}}]{
                \texttt{\tiny{[P1$\langle$GM]}}
                \texttt{OUTPUT: (({-}2, 1), False)} \\
            }
        }
    }
     \\ \\

    \theutterance \stepcounter{utterance}  
    & \multicolumn{4}{p{0.6\linewidth}}{
        \cellcolor[rgb]{0.561,0.737,0.561}{
            \makecell[{{p{\linewidth}}}]{
                \texttt{\tiny{[P1$\rangle$GM]}}
                \texttt{\{"current\_mode":"probing","input\_rationale":"The true outputs at ±1 and false result at {-}2{-}>1 suggest a size{-}based function where all one{-}digit numbers map to 0, such as floor(log10(abs(x))) or digit{-}count minus 1. Testing 10 distinguishes that family.","candidate\_io":\{"input":10,"output":1\},"hypothesis":"def f(x):\textbackslash{}n    \textbackslash{}"\textbackslash{}"\textbackslash{}"Possibly returns the number of decimal digits of abs(x) minus 1, with 0 for abs(x) < 10.\textbackslash{}"\textbackslash{}"\textbackslash{}"\textbackslash{}n    ax = abs(x)\textbackslash{}n    return len(str(ax)) {-} 1","confidence\_score":22\}} \\
            }
        }
    }
    & & \\ \\

    \theutterance \stepcounter{utterance}  
    & & & \multicolumn{4}{p{0.6\linewidth}}{
        \cellcolor[rgb]{0.9,0.9,0.9}{
            \makecell[{{p{\linewidth}}}]{
                \texttt{\tiny{[P1$\langle$GM]}}
                \texttt{OUTPUT: ((99, 3), True)} \\
            }
        }
    }
     \\ \\

    \theutterance \stepcounter{utterance}  
    & \multicolumn{4}{p{0.6\linewidth}}{
        \cellcolor[rgb]{0.561,0.737,0.561}{
            \makecell[{{p{\linewidth}}}]{
                \texttt{\tiny{[P1$\rangle$GM]}}
                \texttt{\{"current\_mode":"probing","input\_rationale":"A strong candidate is: for positive x, return the total number of prime factors counted with multiplicity (Ω); for nonpositive x, return 0. This matches 1{-}>0, {-}1{-}>0, 99=3\^{}2*11 {-}> 3, and excludes {-}2{-}>1. Testing a simple positive prime checks it directly.","candidate\_io":\{"input":2,"output":1\},"hypothesis":"def f(x):\textbackslash{}n    \textbackslash{}"\textbackslash{}"\textbackslash{}"Returns 0 for nonpositive x; otherwise returns the total number of prime factors of x counted with multiplicity.\textbackslash{}"\textbackslash{}"\textbackslash{}"\textbackslash{}n    if x <= 0:\textbackslash{}n        return 0\textbackslash{}n    n = x\textbackslash{}n    c = 0\textbackslash{}n    d = 2\textbackslash{}n    while d * d <= n:\textbackslash{}n        while n \% d == 0:\textbackslash{}n            c += 1\textbackslash{}n            n //= d\textbackslash{}n        d += 1\textbackslash{}n    if n > 1:\textbackslash{}n        c += 1\textbackslash{}n    return c","confidence\_score":63\}} \\
            }
        }
    }
    & & \\ \\

    \theutterance \stepcounter{utterance}  
    & & & \multicolumn{4}{p{0.6\linewidth}}{
        \cellcolor[rgb]{0.9,0.9,0.9}{
            \makecell[{{p{\linewidth}}}]{
                \texttt{\tiny{[P1$\langle$GM]}}
                \texttt{OUTPUT: (({-}1000, 1), False)} \\
            }
        }
    }
     \\ \\

    \theutterance \stepcounter{utterance}  
    & \multicolumn{4}{p{0.6\linewidth}}{
        \cellcolor[rgb]{0.561,0.737,0.561}{
            \makecell[{{p{\linewidth}}}]{
                \texttt{\tiny{[P1$\rangle$GM]}}
                \texttt{\{"current\_mode":"confirming","input\_rationale":"The Ω(n) hypothesis predicts 1000 = 2\^{}3 * 5\^{}3 has 6 prime factors counted with multiplicity. This directly checks the strongest current theory against a previously probed input.","candidate\_io":\{"input":1000,"output":6\},"hypothesis":"def f(x):\textbackslash{}n    \textbackslash{}"\textbackslash{}"\textbackslash{}"Returns 0 for x <= 1, otherwise the total number of prime factors of x counted with multiplicity.\textbackslash{}"\textbackslash{}"\textbackslash{}"\textbackslash{}n    if x <= 1:\textbackslash{}n        return 0\textbackslash{}n    c = 0\textbackslash{}n    d = 2\textbackslash{}n    while d * d <= x:\textbackslash{}n        while x \% d == 0:\textbackslash{}n            c += 1\textbackslash{}n            x //= d\textbackslash{}n        d += 1\textbackslash{}n    if x > 1:\textbackslash{}n        c += 1\textbackslash{}n    return c","confidence\_score":76\}} \\
            }
        }
    }
    & & \\ \\

    \theutterance \stepcounter{utterance}  
    & & & \multicolumn{4}{p{0.6\linewidth}}{
        \cellcolor[rgb]{0.9,0.9,0.9}{
            \makecell[{{p{\linewidth}}}]{
                \texttt{\tiny{[P1$\langle$GM]}}
                \texttt{OUTPUT: (({-}10, 2), True)} \\
            }
        }
    }
     \\ \\

    \theutterance \stepcounter{utterance}  
    & \multicolumn{4}{p{0.6\linewidth}}{
        \cellcolor[rgb]{0.561,0.737,0.561}{
            \makecell[{{p{\linewidth}}}]{
                \texttt{\tiny{[P1$\rangle$GM]}}
                \texttt{\{"current\_mode":"probing","input\_rationale":"A different plausible family is the minimum number of perfect squares needed to sum to |x|, with 0 at ±1. This matches 99{-}>3, {-}10{-}>2, and the negative counterexample at {-}2. Testing a tiny non{-}square like 2 would separate it sharply from prime{-}factor based guesses.","candidate\_io":\{"input":2,"output":2\},"hypothesis":"def f(x):\textbackslash{}n    \textbackslash{}"\textbackslash{}"\textbackslash{}"Possibly returns the minimum number of perfect squares summing to abs(x), with 0 for abs(x) == 1.\textbackslash{}"\textbackslash{}"\textbackslash{}"\textbackslash{}n    n = {-}x if x < 0 else x\textbackslash{}n    if n == 1:\textbackslash{}n        return 0\textbackslash{}n    if n == 0:\textbackslash{}n        return 0\textbackslash{}n    r = int(n ** 0.5)\textbackslash{}n    if r * r == n:\textbackslash{}n        return 0\textbackslash{}n    i = 1\textbackslash{}n    while i * i <= n:\textbackslash{}n        j = int((n {-} i * i) ** 0.5)\textbackslash{}n        if i * i + j * j == n:\textbackslash{}n            return 2\textbackslash{}n        i += 1\textbackslash{}n    m = n\textbackslash{}n    while m \% 4 == 0:\textbackslash{}n        m //= 4\textbackslash{}n    if m \% 8 == 7:\textbackslash{}n        return 4\textbackslash{}n    return 3","confidence\_score":29\}} \\
            }
        }
    }
    & & \\ \\

    \theutterance \stepcounter{utterance}  
    & & & \multicolumn{4}{p{0.6\linewidth}}{
        \cellcolor[rgb]{0.9,0.9,0.9}{
            \makecell[{{p{\linewidth}}}]{
                \texttt{\tiny{[P1$\langle$GM]}}
                \texttt{OUTPUT: ((2, 1), False)} \\
            }
        }
    }
     \\ \\

    \theutterance \stepcounter{utterance}  
    & \multicolumn{4}{p{0.6\linewidth}}{
        \cellcolor[rgb]{0.561,0.737,0.561}{
            \makecell[{{p{\linewidth}}}]{
                \texttt{\tiny{[P1$\rangle$GM]}}
                \texttt{\{"current\_mode":"confirming","input\_rationale":"Two leading explanations fit the revealed data: (1) count prime factors of abs(x) with multiplicity, but return 0 for primes/|x|<=1; or (2) count prime factors of abs(x){-}1. Input 6 separates them cleanly: model (1) gives 2, model (2) gives 1.","candidate\_io":\{"input":6,"output":2\},"hypothesis":"def f(x):\textbackslash{}n    \textbackslash{}"\textbackslash{}"\textbackslash{}"Returns 0 for |x| <= 1 or prime |x|; otherwise returns the number of prime factors of |x| counted with multiplicity.\textbackslash{}"\textbackslash{}"\textbackslash{}"\textbackslash{}n    n = x if x >= 0 else {-}x\textbackslash{}n    if n <= 1:\textbackslash{}n        return 0\textbackslash{}n    m = n\textbackslash{}n    c = 0\textbackslash{}n    d = 2\textbackslash{}n    while d * d <= m:\textbackslash{}n        while m \% d == 0:\textbackslash{}n            c += 1\textbackslash{}n            m //= d\textbackslash{}n        d += 1\textbackslash{}n    if m > 1:\textbackslash{}n        c += 1\textbackslash{}n    return 0 if c == 1 else c","confidence\_score":55\}} \\
            }
        }
    }
    & & \\ \\

    \theutterance \stepcounter{utterance}  
    & & & \multicolumn{4}{p{0.6\linewidth}}{
        \cellcolor[rgb]{0.9,0.9,0.9}{
            \makecell[{{p{\linewidth}}}]{
                \texttt{\tiny{[P1$\langle$GM]}}
                \texttt{OUTPUT: ((0, 0), True)} \\
            }
        }
    }
     \\ \\

    \theutterance \stepcounter{utterance}  
    & \multicolumn{4}{p{0.6\linewidth}}{
        \cellcolor[rgb]{0.561,0.737,0.561}{
            \makecell[{{p{\linewidth}}}]{
                \texttt{\tiny{[P1$\rangle$GM]}}
                \texttt{\{"current\_mode":"confirming","input\_rationale":"The leading hypothesis is: return the total number of prime factors of abs(x) counted with multiplicity, but only for composite numbers; return 0 for {-}1, 0, 1, and primes. A larger prime cleanly tests whether all primes map to 0.","candidate\_io":\{"input":13,"output":0\},"hypothesis":"def f(x):\textbackslash{}n    \textbackslash{}"\textbackslash{}"\textbackslash{}"Returns 0 for abs(x) < 2 or prime abs(x); otherwise returns the total number of prime factors of abs(x) counted with multiplicity.\textbackslash{}"\textbackslash{}"\textbackslash{}"\textbackslash{}n    n = x if x >= 0 else {-}x\textbackslash{}n    if n < 2:\textbackslash{}n        return 0\textbackslash{}n    m = n\textbackslash{}n    c = 0\textbackslash{}n    d = 2\textbackslash{}n    while d * d <= m:\textbackslash{}n        while m \% d == 0:\textbackslash{}n            c += 1\textbackslash{}n            m //= d\textbackslash{}n        d += 1\textbackslash{}n    if m > 1:\textbackslash{}n        c += 1\textbackslash{}n    return 0 if c == 1 else c","confidence\_score":68\}} \\
            }
        }
    }
    & & \\ \\

    \theutterance \stepcounter{utterance}  
    & & & \multicolumn{4}{p{0.6\linewidth}}{
        \cellcolor[rgb]{0.9,0.9,0.9}{
            \makecell[{{p{\linewidth}}}]{
                \texttt{\tiny{[P1$\langle$GM]}}
                \texttt{OUTPUT: ((10, 1), False)} \\
            }
        }
    }
     \\ \\

    \theutterance \stepcounter{utterance}  
    & \multicolumn{4}{p{0.6\linewidth}}{
        \cellcolor[rgb]{0.561,0.737,0.561}{
            \makecell[{{p{\linewidth}}}]{
                \texttt{\tiny{[P1$\rangle$GM]}}
                \texttt{\{"current\_mode":"confirming","input\_rationale":"A strong hypothesis is f(x) = Ω(|x|{-}1), the number of prime factors of |x|{-}1 counted with multiplicity, with 0 when |x|{-}1 < 2. Testing x=4 is decisive because |4|{-}1=3 is prime, so this model predicts 1.","candidate\_io":\{"input":4,"output":1\},"hypothesis":"def f(x):\textbackslash{}n    \textbackslash{}"\textbackslash{}"\textbackslash{}"Returns the number of prime factors of abs(x){-}1 counted with multiplicity, or 0 if abs(x){-}1 < 2.\textbackslash{}"\textbackslash{}"\textbackslash{}"\textbackslash{}n    n = abs(x) {-} 1\textbackslash{}n    if n < 2:\textbackslash{}n        return 0\textbackslash{}n    c = 0\textbackslash{}n    d = 2\textbackslash{}n    while d * d <= n:\textbackslash{}n        while n \% d == 0:\textbackslash{}n            c += 1\textbackslash{}n            n //= d\textbackslash{}n        d += 1\textbackslash{}n    if n > 1:\textbackslash{}n        c += 1\textbackslash{}n    return c","confidence\_score":90\}} \\
            }
        }
    }
    & & \\ \\

    \theutterance \stepcounter{utterance}  
    & & & \multicolumn{4}{p{0.6\linewidth}}{
        \cellcolor[rgb]{0.9,0.9,0.9}{
            \makecell[{{p{\linewidth}}}]{
                \texttt{\tiny{[P1$\langle$GM]}}
                \texttt{OUTPUT: (({-}99, 3), True)} \\
            }
        }
    }
     \\ \\

    \theutterance \stepcounter{utterance}  
    & \multicolumn{4}{p{0.6\linewidth}}{
        \cellcolor[rgb]{0.561,0.737,0.561}{
            \makecell[{{p{\linewidth}}}]{
                \texttt{\tiny{[P1$\rangle$GM]}}
                \texttt{SOLVE: ```python} \\
\texttt{def f(x):} \\
\texttt{    """Returns the number of prime factors of abs(x){-}1, counted with multiplicity; 0 if abs(x){-}1 < 2."""} \\
\texttt{    n = abs(x) {-} 1} \\
\texttt{    if n < 2:} \\
\texttt{        return 0} \\
\texttt{    count = 0} \\
\texttt{    d = 2} \\
\texttt{    while d * d <= n:} \\
\texttt{        while n \% d == 0:} \\
\texttt{            count += 1} \\
\texttt{            n //= d} \\
\texttt{        d += 1} \\
\texttt{    if n > 1:} \\
\texttt{        count += 1} \\
\texttt{    return count} \\
\texttt{```} \\
            }
        }
    }
    & & \\ \\

    \theutterance \stepcounter{utterance}  
    & & & \multicolumn{2}{p{0.3\linewidth}}{
        \cellcolor[rgb]{0.9,0.9,0.9}{
            \makecell[{{p{\linewidth}}}]{
                \texttt{\tiny{[GM$|$GM]}}
                \texttt{Game Verdict: LOSS} \\
            }
        }
    }
    & & \\ \\

    \theutterance \stepcounter{utterance}  
    & & & \multicolumn{2}{p{0.3\linewidth}}{
        \cellcolor[rgb]{0.9,0.9,0.9}{
            \makecell[{{p{\linewidth}}}]{
                \texttt{\tiny{[GM$|$GM]}}
                \texttt{Incorrect.} \\
\texttt{Failed Cases:} \\
\texttt{Input: {[}1000{]} | Expected: 2 | Got: 4} \\
\texttt{Input: {[}{-}2{]} | Expected: 2 | Got: 0} \\
\texttt{Input: {[}{-}1000{]} | Expected: 2 | Got: 4} \\
            }
        }
    }
    & & \\ \\

    \theutterance \stepcounter{utterance}  
    & & & \multicolumn{2}{p{0.3\linewidth}}{
        \cellcolor[rgb]{0.9,0.9,0.9}{
            \makecell[{{p{\linewidth}}}]{
                \texttt{\tiny{[GM$|$GM]}}
                \texttt{def smallest\_divisor\_ge\_two\_abs(x: int) {-}> int:} \\
\texttt{    n = abs(x)} \\
\texttt{    if n < 2:} \\
\texttt{        return 0} \\
\texttt{    for i in range(2, n + 1):} \\
\texttt{        if n \% i == 0:} \\
\texttt{            return i} \\
\texttt{    return 0} \\
            }
        }
    }
    & & \\ \\

    \theutterance \stepcounter{utterance}  
    & & & \multicolumn{2}{p{0.3\linewidth}}{
        \cellcolor[rgb]{0.9,0.9,0.9}{
            \makecell[{{p{\linewidth}}}]{
                \texttt{\tiny{[GM$|$GM]}}
                \texttt{def f(x):} \\
\texttt{    """Returns the number of prime factors of abs(x){-}1, counted with multiplicity; 0 if abs(x){-}1 < 2."""} \\
\texttt{    n = abs(x) {-} 1} \\
\texttt{    if n < 2:} \\
\texttt{        return 0} \\
\texttt{    count = 0} \\
\texttt{    d = 2} \\
\texttt{    while d * d <= n:} \\
\texttt{        while n \% d == 0:} \\
\texttt{            count += 1} \\
\texttt{            n //= d} \\
\texttt{        d += 1} \\
\texttt{    if n > 1:} \\
\texttt{        count += 1} \\
\texttt{    return count} \\
            }
        }
    }
    & & \\ \\

    \theutterance \stepcounter{utterance}  
    & & & \multicolumn{2}{p{0.3\linewidth}}{
        \cellcolor[rgb]{0.9,0.9,0.9}{
            \makecell[{{p{\linewidth}}}]{
                \texttt{\tiny{[GM$|$GM]}}
                \texttt{Mode: passive\_labeled\_pairs} \\
\texttt{Accuracy: 0.260} \\
\texttt{Information requests: 10/15} \\
\texttt{Revealed passive examples: 11} \\
\texttt{Efficiency: 0.400} \\
\texttt{Slack: 5} \\
\texttt{Observed{-}pair consistency: 1.000} \\
\texttt{Membership balance: 6 true / 5 false} \\
\texttt{Parse errors: 0} \\
\texttt{Runtime errors: 0} \\
            }
        }
    }
    & & \\ \\

\end{supertabular}
}

\end{document}
